\documentclass[12pt,a4paper]{article}

\usepackage[utf8]{inputenc}
\usepackage[T1]{fontenc}
\usepackage{amsmath,amssymb,amsthm,mathtools}
\usepackage[colorlinks=true,linkcolor=blue,citecolor=red,urlcolor=blue]{hyperref}
\usepackage{geometry}
\geometry{margin=2.5cm}
\usepackage{enumitem}
\usepackage{booktabs}
\usepackage[capitalise,noabbrev]{cleveref}

\newtheorem{theorem}{Theorem}[section]
\newtheorem{lemma}[theorem]{Lemma}
\newtheorem{proposition}[theorem]{Proposition}
\newtheorem{corollary}[theorem]{Corollary}
\newtheorem{conjecture}[theorem]{Conjecture}
\theoremstyle{definition}
\newtheorem{definition}[theorem]{Definition}
\newtheorem{axiom}[theorem]{Axiom}
\theoremstyle{remark}
\newtheorem{remark}[theorem]{Remark}
\newtheorem{heuristic}[theorem]{Heuristic Principle}

\newcommand{\Ksym}{K_\sigma^{\mathrm{sym}}}
\newcommand{\EE}{\mathbb{E}}
\newcommand{\muG}{\mu_\sigma}
\newcommand{\PP}{\mathcal{P}}

\title{The Riemann Hypothesis:\\
Foundations, Obstructions, and Reorientation}
\author{Lukas Geiger\\
\small Bernau, Germany\\
\small ORCID: \href{https://orcid.org/0009-0005-7296-1534}{0009-0005-7296-1534}}
\date{\today}

\begin{document}
\overfullrule=0pt
\maketitle

\begin{abstract}
This paper, the first in a three-part series on the Riemann Hypothesis, establishes the mathematical landscape and documents the systematic exploration---and ultimately the failure---of a gauge-theoretic approach. In the first part, we develop the analytical framework: the Riemann Hypothesis as a stability problem, the Li criterion as the central tool, and the Bombieri--Lagarias decomposition of the Li coefficients into archimedean and finite-part contributions. Starting from a thermodynamic formalism (Prime Hub Graphs and Ruelle transfer operators), we prove several unconditional structural results: the archimedean--finite decomposition, the archimedean dominance paradox, OS reflection positivity, the arithmetic uncertainty relation, and the spectral census of the arithmetic kernel.

In the second part, we identify four critical obstructions (K1--K4) that individually and collectively doom the original proof strategy: the Gibbs normalization gap, the sign change of the target functional, the wrong monotonicity direction, and the trace-class obstruction. We localize the difficulty precisely through a nine-step logical chain, showing that the \emph{entire} content of RH concentrates in a single step (S5: the unconditional positivity of the Li coefficients). We formulate five concrete open problems and identify the reorientation toward Connes' spectral program as the path forward. Part~II develops this new direction via the Shift Parity Lemma and computer-assisted proofs of even dominance; Part~III draws the final conclusions.
\end{abstract}

\noindent\textbf{MSC 2020:} 11M26, 11M36, 37C30, 47A75, 82B05\\
\textbf{Keywords:} Riemann Hypothesis, Li coefficients, Ruelle transfer operator, archimedean dominance, gauge positivity, obstructions, spectral analysis

\newpage
\tableofcontents
\newpage


% =============================================================================
% =============================================================================
%                     PART I: THE LANDSCAPE
% =============================================================================
% =============================================================================

\part*{The Landscape}
\addcontentsline{toc}{part}{The Landscape}


% =============================================================================
\section{Introduction}
\label{sec:intro}
% =============================================================================

The Riemann Hypothesis (RH) asserts that all non-trivial zeros of the Riemann zeta function $\zeta(s)$ lie on the critical line $\Re(s) = 1/2$. Reformulated via Li's criterion \cite{Li1997}, RH is equivalent to the positivity of an infinite sequence of real numbers:
\begin{equation}
\label{eq:li-criterion}
\text{RH} \quad \Longleftrightarrow \quad \lambda_n \geq 0 \text{ for all } n \geq 1,
\end{equation}
where $\lambda_n = \sum_\rho [1 - (1-1/\rho)^n]$ and the sum runs over the non-trivial zeros of $\zeta$.

This paper is the first in a three-part series:
\begin{itemize}
  \item \textbf{Part~I} (this paper): Foundations, obstructions, and reorientation --- mathematical foundations, analytical tools, structural results, the documentation of four dead ends, the localization of difficulty, and the reorientation toward Connes' spectral program.
  \item \textbf{Part~II} \cite{Geiger2026partII}: Even dominance --- the Shift Parity Lemma, computer-assisted proofs, and the new proof architecture.
  \item \textbf{Part~III} \cite{Geiger2026partIII}: Conclusio --- what is proven, what is excluded, and what remains.
\end{itemize}

The series documents an honest research journey: from a promising gauge-theoretic approach (Sections~\ref{sec:thermodynamic}--\ref{sec:gauge}), through its failure at multiple points (Sections~\ref{sec:K1}--\ref{sec:K4}), to a new structural understanding of the problem via the Weil quadratic form (Section~\ref{sec:reorientation}), and a final accounting of what has been achieved (Part~III).

\subsection{The Hope}

The Hilbert--P\'olya conjecture proposes that the non-trivial zeros correspond to eigenvalues of a self-adjoint operator. The Berry--Keating semiclassical approach ($\hat{H} = xp$) fails due to the Weyl law obstruction. This paper shifts the paradigm from quantum mechanics to statistical mechanics: the zeros are generated not as eigenvalues of a Hamiltonian, but as resonances of a Ruelle-type transfer operator.

The guiding idea is that the Riemann Hypothesis is a \emph{stability condition}: the arithmetic structure of the primes, encoded in the Li coefficients $\lambda_n$, is stabilized by a competition between arithmetic instability (from the primes) and archimedean convexity (from the Gamma function). Understanding this competition is the central theme of the series.


% =============================================================================
\section{The Thermodynamic Formalism}
\label{sec:thermodynamic}
% =============================================================================

\subsection{Prime Hub Graph and the Vacuum Trace}

To reproduce the Euler product $\zeta(s) = \prod_p (1-p^{-s})^{-1}$, we construct a dynamical system whose primitive periodic orbits correspond to primes.

\begin{definition}[Prime Orbit Axiom]
We postulate a topologically mixing directed graph $\mathcal{G} = (\mathcal{V}, \mathcal{E})$ whose primitive cycles $\gamma_p$ are in bijection with the primes $p \in \PP$, with geometric length $\ell(\gamma_p) = \log p$.
\end{definition}

\begin{remark}[Motivation vs.\ logical dependency]
The Prime Orbit Axiom is a constructive motivation. The analytical results in subsequent sections depend solely on properties of $\zeta(s)$ and the arithmetic Euler product, and are logically independent of this graph construction.
\end{remark}

Using a holonomy filter on the free monoid generated by $\PP$, we define a transfer operator $\mathcal{L}_s$ and prove:

\begin{theorem}[Recovery of the Zeta Function]
\label{thm:vacuum-trace}
For $\Re(s) > 1$, the Fredholm determinant of the vacuum-projected transfer operator satisfies:
\[
\det(I - \mathcal{L}_s)^{\mathrm{vac}} = \prod_p (1 - p^{-s}) = \frac{1}{\zeta(s)}.
\]
\end{theorem}

\begin{proof}
By the Ruelle identity \cite{Ruelle1976}, $\log\det(I - \mathcal{L}_s) = -\sum_{m=1}^\infty \frac{1}{m}\operatorname{Tr}(\mathcal{L}_s^m)$. The vacuum projection restricts the trace to configurations that form closed loops of prime powers: the only periodic orbits of length $m\log p$ in $\mathcal{G}$ are the $m$-fold traversals of the primitive cycle $\gamma_p$, each contributing weight $p^{-ms}$ with multiplicity $\log p$ (the geometric length). Summing over prime powers and exponentiating:
\[
  \log\det(I - \mathcal{L}_s)^{\mathrm{vac}} = -\sum_p \sum_{m=1}^\infty \frac{1}{m} p^{-ms} = \sum_p \log(1 - p^{-s}).
\]
This is the logarithm of the Euler product $\prod_p(1-p^{-s}) = 1/\zeta(s)$, valid for $\Re(s) > 1$ by absolute convergence.
\end{proof}

\subsection{The MEPP Heuristic}

The Maximum Entropy Production Principle (MEPP) interprets the critical line $\Re(s) = 1/2$ as a detailed-balance condition:

\begin{heuristic}[Model RH via MEPP]
If the primes represent the optimally stable configuration of a ``prime gas'' subject to the PNT density constraint, the resulting Gibbs measure satisfies time-reversal symmetry, forcing zeros onto the critical line.
\end{heuristic}

This heuristic does not constitute a proof but motivates the thermodynamic perspective developed below.


% =============================================================================
\section{The Li Coefficients: Analytical Framework}
\label{sec:li-coefficients}
% =============================================================================

\subsection{Bombieri--Lagarias Representation}

The completed xi function $\xi(s) = \frac{1}{2}s(s-1)\pi^{-s/2}\Gamma(s/2)\zeta(s)$ provides the definitive representation:

\begin{theorem}[Bombieri--Lagarias \cite{BombieriLagarias1999}]
\label{thm:bombieri-lagarias}
For $n \geq 1$:
\[
\lambda_n = \frac{1}{(n-1)!}\,\frac{\partial^n}{\partial s^n}\Big[s^{n-1}\log\xi(s)\Big]_{s=1}.
\]
Since $\xi(s)$ is entire of order~$1$ with $\xi(1) = 1/2 \neq 0$, $\log\xi(s)$ is holomorphic near $s = 1$ and the derivative is well-defined.
\end{theorem}

\subsection{Archimedean--Finite Decomposition}

\begin{proposition}[Archimedean--Finite Decomposition]
\label{prop:decomposition}
Each Li coefficient decomposes as $\lambda_n = \lambda_n^{(\infty)} + \lambda_n^{(\mathrm{fin})}$, where:
\begin{align}
\lambda_n^{(\infty)} &= \frac{1}{(n-1)!}\frac{\partial^n}{\partial s^n}\Big[s^{n-1}\log\!\big(\tfrac{1}{2}s\,\pi^{-s/2}\,\Gamma(s/2)\big)\Big]_{s=1}, \label{eq:arch-part} \\
\lambda_n^{(\mathrm{fin})} &= \frac{1}{(n-1)!}\frac{\partial^n}{\partial s^n}\Big[s^{n-1}\log g(s)\Big]_{s=1}, \quad g(s) := (s-1)\zeta(s). \label{eq:fin-part}
\end{align}
Both parts are individually computable; $g(s)$ is holomorphic at $s=1$ with $g(1) = 1$.
\end{proposition}

\begin{remark}[Three-regime structure]
\label{rem:three-regime}
Numerical computation reveals:
\begin{enumerate}[label=(\alph*)]
  \item \textbf{Small $n$ ($n \lesssim 8$):} $\lambda_n^{(\infty)} < 0$, compensated by $\lambda_n^{(\mathrm{fin})} > 0$. Positivity is a delicate cancellation.
  \item \textbf{Transition ($n \approx 8$):} $\lambda_n^{(\infty)}$ crosses zero.
  \item \textbf{Large $n$ ($n \gg 10$):} $\lambda_n^{(\infty)} \sim \frac{n}{2}\log n$ dominates; positivity is automatic.
\end{enumerate}
The ``hard'' instances of Li's criterion concentrate in regime~(a). Table~\ref{tab:li-coefficients} shows the decomposition for $n = 1, \ldots, 13$.
\end{remark}

\begin{table}[h]
\centering
\renewcommand{\arraystretch}{1.15}
\begin{tabular}{r|r|r|r|r}
$n$ & $\lambda_n$ & $\lambda_n^{(\infty)}$ & $\lambda_n^{(\mathrm{fin})}$ & $c_n = \lambda_n/n$ \\\hline
1 & $0.023$ & $-0.554$ & $0.577$ & $0.023$ \\
2 & $0.092$ & $-0.875$ & $0.967$ & $0.046$ \\
3 & $0.208$ & $-1.013$ & $1.221$ & $0.069$ \\
5 & $0.576$ & $-0.883$ & $1.458$ & $0.115$ \\
8 & $1.466$ & $+0.021$ & $1.445$ & $0.183$ \\
10 & $2.279$ & $+0.956$ & $1.324$ & $0.228$ \\
13 & $3.817$ & $+2.725$ & $1.093$ & $0.294$ \\
\end{tabular}
\caption{Archimedean--finite decomposition of the Li coefficients. The sequence $c_n = \lambda_n/n$ is strictly increasing, ruling out complete monotonicity.}
\label{tab:li-coefficients}
\end{table}


% =============================================================================
\section{The Archimedean Dominance Paradox}
\label{sec:dominance}
% =============================================================================

\begin{proposition}[Stability Decomposition]
\label{prop:stability}
The global curvature $\mathcal{M}(\sigma) = \partial_s^2\log\xi(s)$ decomposes as
$\mathcal{M} = \mathcal{M}_{\mathrm{arith}} + \mathcal{M}_{\mathrm{arch}}$:
\begin{enumerate}
  \item \textbf{Arithmetic instability:} $\mathcal{M}_{\mathrm{arith}}(1) = \sum_\rho \Re(1/\rho^2) < 0$.
  \item \textbf{Archimedean stability:} $\mathcal{M}_{\mathrm{arch}}(\sigma) > 0$ for $\sigma > 1/2$ (explicitly positive and convex).
\end{enumerate}
\end{proposition}

\noindent The Riemann Hypothesis is equivalent to the archimedean convexity globally dominating the arithmetic concavity for all $\sigma > 1/2$. This ``dominance paradox'' --- a manifestly positive continuous term must overcome infinitely many oscillatory arithmetic contributions --- is the structural core of the problem.


% =============================================================================
\section{Structural Results}
\label{sec:structural}
% =============================================================================

We establish five unconditional results that constrain the problem and form the toolkit for subsequent sections and papers.

\subsection{Prime-Sum Functionals and the Gibbs Framework}

\begin{definition}[Arithmetic Gibbs Measure]
The Gibbs measure is $\muG(p) = p^{-\sigma}/P(\sigma)$ with $P(\sigma) = \sum_p p^{-\sigma}$. The prime-sum functionals are:
\begin{align*}
A(\sigma) &:= \sum_p \frac{\log p}{p^\sigma(p^\sigma - 1)}, &
B(\sigma) &:= \sum_p \frac{\log p}{p^\sigma},\\
C(\sigma) &:= \sum_p \frac{(\log p)^2}{p^\sigma(p^\sigma - 1)}, &
D(\sigma) &:= \sum_p \frac{(\log p)^2}{(p^\sigma - 1)^2}.
\end{align*}
\end{definition}

\begin{lemma}[Finite-Part Convergence]
\label{lem:convergence}
$A(\sigma)$, $C(\sigma)$, and $D(\sigma)$ converge absolutely at $\sigma = 1$ and are right-continuous. $B(\sigma)$ and $P(\sigma)$ diverge as $\sim\log(1/(\sigma-1))$.
\end{lemma}

\begin{remark}[Gibbs normalization gap]
The Gibbs-normalized expectation $-A(\sigma)/P(\sigma)$ loses the finite-part information as $\sigma \to 1^+$: the $1/P(\sigma)$ normalization annihilates $\lambda_n^{(\mathrm{fin})}$. This is the first hint that the thermodynamic approach has structural limitations (see Section~\ref{sec:K1}).
\end{remark}

\subsection{Non-Identification Theorem}

\begin{proposition}[Non-Identification]
\label{prop:non-ident}
The excess free energy $I_n := \lim_{\sigma \to 1^+}\Phi_n(\sigma)$ and the Li coefficient $\lambda_n$ are generically \emph{not} equal: $I_n \neq \lambda_n$. The former is a KL divergence (tilt cost), the latter a spectral trace of $\log\xi$.
\end{proposition}

\subsection{Variance Lower Bound}

\begin{proposition}[Variance Bound]
\label{prop:variance}
$I_1 \geq \frac{1}{2}V_1$ where $V_1 = \lim_{\sigma \to 1^+}\mathrm{Var}_{\muG}(H_{1,\sigma}) > 0$. Numerically, $I_1 \geq 0.034 > \lambda_1 \approx 0.023$.
\end{proposition}

\subsection{OS Reflection Positivity}

\begin{proposition}[OS Reflection Positivity]
\label{prop:os}
The Osterwalder--Schrader form
\[
K(t,u) = \sum_{p \in \PP} w_p\,e^{-i(u-t)\log p}, \qquad w_p = \frac{p^{-(\sigma+1)}}{P(\sigma)} > 0,
\]
is positive definite on $L^2_+(\mathbb{R})$. The OS Hilbert space carries a contractive semigroup $e^{-\tau H}$ with $H = \mathrm{diag}(\log p)$.
\end{proposition}

\begin{proof}
By Bochner's theorem: $K(t,u) = \hat{\mu}(u-t)$ with $\mu = \sum_p w_p\,\delta_{\log p}$ a positive measure, so $\int\!\!\int \bar{f}(t)f(u)K(t,u)\,dt\,du = \sum_p w_p|\hat{f}(\log p)|^2 \geq 0$.
\end{proof}

\subsection{Arithmetic Uncertainty Relation}

\begin{heuristic}[Arithmetic Uncertainty Relation]
\label{heur:uncertainty}
For normalized $\psi \in \ell^2(\PP)$, define the prime-domain spread $\Delta_D$ via
\[
\Delta_D^2 = \sum_p |a_p|^2(\log p - \overline{\log p})^2
\]
and the zero-domain spread $\Delta_t$ via the Fourier transform $\hat\psi(t) = \sum_p a_p e^{-it\log p}$. Then $\Delta_D \cdot \Delta_t \geq 1/2$.
\end{heuristic}

\begin{remark}
This follows formally from the substitution $x_p = \log p$ and the Heisenberg--Kennard--Weyl inequality on $L^2(\mathbb{R})$. However, $\psi$ is supported on the discrete set $\{\log p : p \in \PP\}$, not on all of~$\mathbb{R}$, so the reduction requires an embedding into $L^2(\mathbb{R})$ whose properties depend on the distribution of primes. We state this as a heuristic principle, not a theorem.

The physical content is a consistency check: a zero off the critical line would require a ``localized phase conspiracy'' among finitely many primes, which the arithmetic incommensurability of $\{\log p\}$ makes implausible but does not rigorously exclude.
\end{remark}

\subsection{Strict Convexity and the Spectral Census}

\begin{proposition}[Strict Convexity]
\label{prop:convexity}
$F(\sigma) := -\log\zeta(\sigma)$ is strictly convex on $(1,\infty)$: $F''(\sigma) = \sum_p (\log p)^2 p^{-\sigma}/(1-p^{-\sigma})^2 > 0$.
\end{proposition}

\begin{proposition}[Spectral Census]
\label{prop:spectral-census}
For the symmetrized arithmetic kernel $\Ksym$ truncated to $N \geq 300$ primes and $\sigma \in (1, 5)$, the negative inertia index is exactly~$2$, independent of $N$ and $\sigma$.
\end{proposition}


% =============================================================================
\section{The Gauge Equivalence Framework}
\label{sec:gauge}
% =============================================================================

The spectral census (\Cref{prop:spectral-census}) suggests a gauge-correction approach: can a rank-2 correction restore positive-semidefiniteness?

\begin{definition}[Admissible Gauge]
A symmetric kernel $G_\sigma$ on $\PP \times \PP$ is admissible if $\sum_{p,q}G_\sigma(p,q)\muG(p)\muG(q) = 0$.
\end{definition}

\begin{theorem}[Gauge Equivalence]
\label{thm:gauge}
For fixed $\sigma > 1$, the following are equivalent:
\begin{enumerate}[label=(\roman*)]
  \item The target functional $F(\sigma) = A(\sigma)B(\sigma) - P(\sigma)[C(\sigma)+D(\sigma)] \geq 0$.
  \item There exists an admissible gauge making $\Ksym + G_\sigma \succeq 0$.
\end{enumerate}
\end{theorem}

This tautological equivalence provides the formal structure for the gauge-theoretic approach. Whether it leads anywhere depends on the \emph{global} behavior of $F(\sigma)$ --- a question addressed in Section~\ref{sec:K2}, where we discover that $F(\sigma) < 0$ for $\sigma > 1.14$.


% =============================================================================
\section{Summary: The Starting Position}
\label{sec:summary}
% =============================================================================

At the end of the landscape analysis, we have assembled the following toolkit:

\begin{center}
\renewcommand{\arraystretch}{1.15}
\begin{tabular}{ll}
\toprule
\textbf{Result} & \textbf{Status} \\
\midrule
R1: Bombieri--Lagarias representation & proven \\
R2: Archimedean--finite decomposition & proven \\
R3: Non-identification ($I_n \neq \lambda_n$) & proven \\
R4: Variance lower bound ($I_1 \geq \frac{1}{2}V_1$) & proven \\
R5: Arithmetic uncertainty relation & heuristic \\
R6: Strict convexity of $-\log\zeta$ & proven \\
R7: OS reflection positivity & proven \\
R8: Spectral census (index $= 2$) & numerical \\
R9: Gauge equivalence (tautological) & proven \\
\bottomrule
\end{tabular}
\end{center}

The hope is clear: use the gauge framework (R9) together with the spectral census (R8) to construct a global certificate proving $\lambda_n \geq 0$. The following sections will show that this hope is dashed by a sign change at $\sigma_0 \approx 1.14$ --- but that a completely different approach, through Connes' Weil quadratic form, opens a new and more promising path.


% =============================================================================
% =============================================================================
%                     PART II: THE DEAD ENDS
% =============================================================================
% =============================================================================

\part*{The Dead Ends}
\addcontentsline{toc}{part}{The Dead Ends}


% =============================================================================
\section{The Original Strategy}
\label{sec:original-strategy}
% =============================================================================

The toolkit of nine structural results (R1--R9) and the gauge equivalence framework established above suggest a clear proof strategy:

\begin{enumerate}
  \item Use the spectral census (R8: negative inertia index $= 2$) to identify the unstable modes.
  \item Construct an explicit rank-2 gauge correction restoring positive-semidefiniteness.
  \item Apply the gauge equivalence theorem (R9) to deduce $\lambda_n \geq 0$.
\end{enumerate}

This strategy fails at step~2. In this part, we document four independent obstructions that kill it, then use the resulting understanding to localize the difficulty precisely.


% =============================================================================
\section{Dead End 1: The Gibbs Normalization Gap (K1)}
\label{sec:K1}
% =============================================================================

The thermodynamic formulation defines the renormalized constraint $m_n(\sigma) = \EE_{\muG}[X_{n,\sigma}] - h_n(\sigma)$, with the hope that $\lim_{\sigma \to 1^+} m_n(\sigma) = \lambda_n$.

\begin{proposition}[K1: Critical Identification Fails]
\label{prop:K1}
The Gibbs-normalized expectation $m_1(\sigma) = -A(\sigma)/P(\sigma) - h_1(\sigma)$ does \emph{not} converge to $\lambda_1$ as $\sigma \to 1^+$. The $1/P(\sigma)$-normalization annihilates the finite-part contribution $\lambda_1^{(\mathrm{fin})} = \gamma \approx 0.577$.
\end{proposition}

\begin{proof}
The numerator sum $S(\sigma) = \sum_p (\log p)p^{-\sigma}f(p,\sigma)$ converges to a finite limit, while $P(\sigma) \to \infty$. Therefore $-S(\sigma)/P(\sigma) \to 0$. The Li coefficient $\lambda_1$ arises from $\frac{d}{ds}\log\xi(s)|_{s=1}$, which involves $\zeta'(s)/\zeta(s) + 1/(s-1) \to \gamma$ --- a cancellation that occurs \emph{before} Gibbs normalization, not after.
\end{proof}

\begin{remark}[Consequence]
The excess free energy $I_n$ and the Li coefficient $\lambda_n$ are structurally distinct objects (\Cref{prop:non-ident}). Even $I_n \geq 0$ (which holds unconditionally as a KL divergence) does not imply $\lambda_n \geq 0$. This is the ``convexity trap'': by the Bregman inequality, $I_n \geq 0$ is a universal identity for any strictly convex generating function, independent of~$\mathrm{sign}(\lambda_n)$.
\end{remark}


% =============================================================================
\section{Dead End 2: Global Gauge Positivity Fails (K2)}
\label{sec:K2}
% =============================================================================

The gauge equivalence theorem (R9) requires $F(\sigma) \geq 0$ for all $\sigma > 1$.

\begin{proposition}[K2: Sign Change of $F(\sigma)$]
\label{prop:K2}
The target functional $F(\sigma) = A(\sigma)B(\sigma) - P(\sigma)[C(\sigma) + D(\sigma)]$ satisfies:
\begin{itemize}
  \item $F(\sigma) > 0$ for $\sigma \in (1, \sigma_0)$ with $\sigma_0 \approx 1.14$,
  \item $F(\sigma) < 0$ for all tested $\sigma > \sigma_0$ (minimum $F \approx -0.184$ at $\sigma \approx 1.28$).
\end{itemize}
This sign change is independent of RH.
\end{proposition}

\begin{proof}
Near $\sigma = 1 + \varepsilon$: $A \cdot B \sim 1/\varepsilon$ dominates $P(C+D) \sim \log(1/\varepsilon)$, so $F > 0$. For larger $\sigma$: $B(\sigma)$ decays faster than $P(\sigma)$, and the product $AB$ falls below $P(C+D)$. Numerical computation with $N = 10000$ primes and Rosser--Schoenfeld tail bounds confirms the transition at $\sigma_0 \approx 1.14$.
\end{proof}

\begin{remark}[Consequence]
The claim ``RH $\Leftrightarrow$ $F(\sigma) \geq 0$ for all $\sigma > 1$'' is \textbf{false}. The gauge approach works only in the boundary layer $(1, 1.14)$, not globally. No admissible gauge can make $\Ksym$ positive-semidefinite for $\sigma > 1.14$.
\end{remark}


% =============================================================================
\section{Dead End 3: Monotonicity Direction Wrong (K3)}
\label{sec:K3}
% =============================================================================

The original argument assumed that $\frac{d}{ds}\log\xi(\sigma)$ is decreasing, making $\lambda_1$ the maximum.

\begin{proposition}[K3: Convexity, Not Concavity]
\label{prop:K3}
$\frac{d^2}{ds^2}\log\xi(\sigma) \approx +0.046 > 0$ at $\sigma = 1$. The function $\log\xi$ is \emph{convex} near $s = 1$, not concave. Therefore $\lambda_1$ is the \emph{minimum} of the derivative profile, not the maximum.
\end{proposition}

\begin{remark}[Consequence]
This reverses the logic of the monotonicity argument: instead of showing $\lambda_1 \geq 0$ by bounding the maximum from below, one would need to bound the minimum from below --- a strictly harder problem.
\end{remark}


% =============================================================================
\section{Dead End 4: The Stieltjes Realization and A8 (K4)}
\label{sec:K4}
% =============================================================================

\subsection{Path D2: Stieltjes Realization}

\begin{proposition}[K4a: Complete Monotonicity Fails]
The sequence $c_n = \lambda_n/n$ is \emph{strictly increasing}: $c_1 \approx 0.023$, $c_{13} \approx 0.294$, with $c_n \sim \frac{1}{2}\log n$ asymptotically. This violates complete monotonicity at $k = 1$ ($\Delta c_n > 0$, so $(-1)^1\Delta c_n < 0$). The Stieltjes realization approach with a measure on~$[0,1]$ is ruled out.
\end{proposition}

Under RH, the spectral measure lives on the unit circle $|w_\rho| = 1$, not $[0,1]$. The positivity of $c_n$ is an essentially number-theoretic property.

\subsection{K4b: Inconsistency of the Original A8}

\begin{proposition}[Hardy Contradiction]
The original axiom A8 posits $\mathcal{K}_0 \geq 0$ trace-class with $\xi(1/2+it) = C\cdot\det(I + V_t\mathcal{K}_0V_t^*)$. But $\mathcal{K}_0 \geq 0$ implies $\det(I + V_t\mathcal{K}_0V_t^*) \geq 1 > 0$ for all real $t$. This contradicts Hardy's theorem (1914): $\xi$ has infinitely many zeros on the critical line.
\end{proposition}

\subsection{The Trace-Class Obstruction}

The corrected axiom A8$'$ requires $\xi(s)/\xi(0) = \det(I - \mathcal{L}_s)$ with $\mathcal{L}_s$ nuclear. However, the Euler product gives $\mathcal{L}_s = \mathrm{diag}(p^{-s})$ on $\ell^2(\PP)$, which is trace-class only for $\Re(s) > 1$. At the critical line, $\sum_p p^{-1/2} = +\infty$.

For the Selberg zeta function, the analogous construction works because geodesic lengths grow exponentially ($\#\{\gamma: l_\gamma \leq T\} \sim e^T/T$), while prime ``lengths'' $\log p$ are too dense ($\pi(x) \sim x/\log x$). This density gap is the fundamental obstruction.


% =============================================================================
\section{Route Analysis: What Survives}
\label{sec:routes}
% =============================================================================

The axiom system (defined in Part~II) admits two logical routes to RH:

\begin{description}[font=\bfseries,leftmargin=2em]
  \item[Route 1 (A7 + A8$'$):] A7 (OS reflection positivity) is \textbf{proven}. A8$'$ (nuclear Fredholm determinant for $\xi$) remains \textbf{open} and faces the trace-class obstruction.
  \item[Route 2 (A9):] \textbf{Dead.} A9(i--ii) hold, but A9(iii) is broken by K1 and the convexity trap.
\end{description}


% =============================================================================
\section{The Localization Table}
\label{sec:localization}
% =============================================================================

We trace the logical chain from definitions to RH, identifying exactly where the difficulty concentrates.

\begin{center}
\renewcommand{\arraystretch}{1.2}
\begin{tabular}{c|l|l|l}
\textbf{Step} & \textbf{Statement} & \textbf{Status} & \textbf{Method} \\\hline
S1 & Define $\xi(s)$ & proven & definition \\
S2 & $\lambda_n$ via Bombieri--Lagarias & proven & \cite{BombieriLagarias1999} \\
S3 & $\lambda_n = \lambda_n^{(\infty)} + \lambda_n^{(\mathrm{fin})}$ & proven & Section~\ref{sec:li-coefficients} \\
S4 & $\lambda_n^{(\infty)} \sim \frac{n}{2}\log n$ for large $n$ & proven & Stirling \\
\textbf{S5} & \textbf{$\lambda_n \geq 0$ for all $n \geq 1$} & \textbf{OPEN} & \textbf{$\Longleftrightarrow$ RH} \\
S6 & All zeros on $\Re(s) = 1/2$ & $\Leftrightarrow$ S5 & Li (1997) \\
S7 & $\pi(x) = \mathrm{Li}(x) + O(\sqrt{x}\log x)$ & $\Leftarrow$ S6 & von Koch \\
S8 & Spectral census (index $= 2$) & numerical & Section~\ref{sec:structural} \\
S9 & Gauge equivalence (tautological) & proven & Section~\ref{sec:gauge} \\
\end{tabular}
\end{center}

\noindent\textbf{Conclusion:} The \emph{entire} content of the Riemann Hypothesis concentrates in Step S5. Everything before it is unconditional; everything after it is either equivalent or a consequence. The present program has not reduced S5 to a simpler statement.


% =============================================================================
\section{Five Open Problems}
\label{sec:open-problems}
% =============================================================================

Based on our analysis, we formulate five concrete open problems ordered by estimated difficulty:

\begin{conjecture}[OP1: Analytical Proof of Index Rigidity]
The negative inertia index of $\Ksym$ is exactly~$2$ for all $N$ sufficiently large and $\sigma \in (1, \sigma_{\max})$. A proof should exploit the rank-2 hub structure of the transfer operator.
\end{conjecture}

\begin{conjecture}[OP2: Uniform Bound on $\lambda_n^{(\mathrm{fin})}$]
$\lambda_n^{(\mathrm{fin})} > 0$ for all $n \geq 1$. This is supported by the numerical observation that $\lambda_n^{(\mathrm{fin})}$ is always positive with slow algebraic decay.
\end{conjecture}

\begin{conjecture}[OP3: Growth Rate]
$\lambda_n^{(\infty)} \sim \frac{n}{2}\log n$ with explicit error bounds, uniform in $n$. This would settle RH for all $n$ above an explicit threshold.
\end{conjecture}

\begin{conjecture}[OP4: Small-$n$ Positivity]
$\lambda_n > 0$ for $n = 1, \ldots, n_0$ (some explicit $n_0$) unconditionally --- without assuming RH. The first $10^4$ coefficients are known to be positive \cite{Keiper1992}.
\end{conjecture}

\begin{conjecture}[OP5: Nuclear Transfer Operator for $\xi$]
Construct a separable Hilbert space $\mathcal{V}$ and a nuclear family $s \mapsto \mathcal{L}_s$ with $\xi(s)/\xi(0) = \det(I - \mathcal{L}_s)$, spectral radius $< 1$ for $\Re(s) > 1/2$, and eigenvalue $1$ at the zeros.
\end{conjecture}

OP5 is essentially the Hilbert--P\'olya program and connects to Deninger's foliated spaces \cite{Deninger2002} and Connes' adelic approach \cite{Connes1999}.


% =============================================================================
\section{Reorientation: Connes' Spectral Program}
\label{sec:reorientation}
% =============================================================================

The failure of the gauge-theoretic approach (K1--K4) forces a fundamental reorientation. A recent survey by Connes \cite{Connes2026} provides a new path: the Weil quadratic form $\mathrm{QW}_\lambda$ on $L^2([-\log\lambda, \log\lambda])$, constructed from the Weil explicit formula, defines a self-adjoint operator $A_\lambda$ whose spectral properties are directly linked to the zeros of $\zeta$.

\begin{theorem}[Connes, Theorem 6.1 in \cite{Connes2026}]
If the minimum eigenvalue of $A_\lambda$ is simple with even eigenfunction, then the Fourier transform of the eigenfunction has all zeros on the real line.
\end{theorem}

Connes' Section~6.6 identifies the verification of this \emph{even dominance} property as the remaining open step. Crucially, Connes' Hurwitz argument requires even dominance only along a sequence $\lambda_n \to \infty$, not for all $\lambda$.

This is where the thermodynamic insight pays off in an unexpected way: the decomposition analysis from the landscape sections (the prime-sum functionals, the trigonometric structure of shift operators) becomes the key to understanding \emph{why} even dominance holds. Part~II develops this completely, establishing the Shift Parity Lemma, the frontier-prime mechanism, and computer-assisted proofs of even dominance for $33$ values of $\lambda$ up to $1{,}300{,}000$.

\begin{remark}[From despair to hope]
The gauge-theoretic approach failed because it tried to prove a \emph{global} statement ($\lambda_n \geq 0$ for all $n$) via a \emph{local} tool (the Gibbs measure near $\sigma = 1$). The Connes approach succeeds where the gauge approach fails because it works directly in the spectral domain of $A_\lambda$, where the even-odd decomposition provides a natural structure. The Shift Parity Lemma (Part~II) shows that this structure is fundamentally algebraic, not analytic.
\end{remark}


% =============================================================================
\section{Conclusion}
\label{sec:conclusion}
% =============================================================================

This paper has documented both a toolkit and an honest failure. The gauge-theoretic approach to the Riemann Hypothesis produced genuine structural insights (R1--R9) but cannot bridge the gap between the local thermodynamic analysis and the global spectral positivity required by Li's criterion. Four independent obstructions (K1--K4) kill the original strategy, and the localization table shows that no reduction to a simpler problem has been achieved: the difficulty remains concentrated in Step~S5.

However, the failure is productive:
\begin{enumerate}
  \item The landscape analysis (Sections~\ref{sec:thermodynamic}--\ref{sec:gauge}) establishes a permanent toolkit of unconditional structural results.
  \item The dead ends (K1--K4) are clearly mapped, preventing others from pursuing them.
  \item The reorientation to Connes' program leverages the thermodynamic insight (prime shift operators, trigonometric structure) in a new context.
  \item The understanding that even dominance is driven by prime contributions, not archimedean effects, emerged directly from the decomposition analysis.
\end{enumerate}

Part~II takes up the new thread: the Shift Parity Lemma, the computer-assisted proof of even dominance, and the identification of the cumulative step as the remaining challenge. Part~III draws the final conclusions on what has been proven, what has been excluded, and what remains open.


\section*{AI Disclosure}

This manuscript was developed in extensive collaboration with AI language models (Claude, Anthropic; Copilot, Microsoft/GitHub; Gemini, Google DeepMind). AI contributed to conceptual exploration, analytical work, literature review, writing, and critical review. The author, Lukas Geiger, conceived the research direction, provided domain expertise, and exercised final editorial authority. The author assumes full and sole scholarly responsibility for all content, including any errors.


\begin{thebibliography}{99}

\bibitem{Li1997}
X.-J.~Li, \emph{The positivity of a sequence of numbers and the Riemann hypothesis}, J.~Number Theory~\textbf{65} (1997), 325--333.

\bibitem{BombieriLagarias1999}
E.~Bombieri and J.~C.~Lagarias, \emph{Complements to Li's criterion for the Riemann hypothesis}, J.~Number Theory~\textbf{77} (1999), 274--287.

\bibitem{Ruelle1976}
D.~Ruelle, \emph{Zeta-functions for expanding maps and Anosov flows}, Invent.~Math.~\textbf{34} (1976), 231--242.

\bibitem{Keiper1992}
J.~B.~Keiper, \emph{Power series expansions of Riemann's $\xi$ function}, Math.~Comp.~\textbf{58} (1992), 765--773.

\bibitem{Connes1999}
A.~Connes, \emph{Trace formula in noncommutative geometry and the zeros of the Riemann zeta function}, Selecta Math.~\textbf{5} (1999), 29--106.

\bibitem{Connes2026}
A.~Connes, \emph{The Riemann Hypothesis: Past, Present and a Letter Through Time}, arXiv:2602.04022v1, 2026.

\bibitem{Deninger2002}
C.~Deninger, \emph{On the nature of the ``explicit formulas'' in analytic number theory}, in: Number Theoretic Methods, Developments in Mathematics, vol.~\textbf{8}, Kluwer Academic Publishers (2002), 97--118.

\bibitem{Geiger2026partII}
L.~Geiger, \emph{The Riemann Hypothesis: Even Dominance of the Weil Quadratic Form}, 2026.

\bibitem{Geiger2026partIII}
L.~Geiger, \emph{The Riemann Hypothesis: Conclusio}, 2026.

\end{thebibliography}

\end{document}
